\documentclass[italian,12pt]{article}

\def\Title{Verbale interno}
\def\Subtitle{Riunione interna settimanale}
\def\Author{7Last}
\def\Date{2024-05-22}
\def\Version{v1.0}

%pacchetti extra da scaricare dblfloatfix, fancyhdr
\usepackage[left=2cm, right=2cm, bottom=3cm, top=3cm]{geometry}
\usepackage{fancyhdr}%creazione header-footer
\usepackage{graphicx} %serve per inserire immagini
\graphicspath{ {../../../logo/} }
%\usepackage{dblfloatfix} %serve per posizionare gli elementi dove si vuole
\usepackage[hidelinks]{hyperref} %serve per i link
\usepackage{tikz}
\usepackage{tgadventor} % font
\usepackage[useregional=numeric,showseconds=true,showzone=false]{datetime2}
\usepackage{caption}

\usepackage{hyperref}
\usepackage{tocloft}
\usepackage{titlesec}
\usepackage{color}
\usepackage{ulem}
\usepackage{pgfplots}
\usepackage{pgf-pie}
\usepackage[italian]{babel}
\usepackage{comment}
\usepackage{tabularx}
\usepackage{longtable}
\usepackage{float}
\usepackage{amsmath}

\linespread{1.2}
\captionsetup[table]{labelformat=empty}
\geometry{headsep=1.5cm}

\renewcommand{\contentsname}{Indice}
\renewcommand\familydefault{\sfdefault}

\begin{document}

\newgeometry{left=2cm,right=2cm,bottom=2.1cm,top=2.1cm}
\begin{titlepage}
    \vspace*{.5cm}

    \vspace{2cm}
    {
        \centering
        {\bfseries\huge \Title\par}
        \bigbreak
        {\bfseries\large \Author\par}
        \bigbreak
        {\Version\par}
        \vfill

        \begin{tikzpicture}[remember picture,overlay]

            \fill[blue!20!black] (current page.south west) -- ++(10cm,0) -- ++(-10cm,15cm) -- cycle;
            \fill[orange] (current page.south east) -- ++(-18cm,0) -- ++(21.6cm,22cm) -- cycle;

            \clip (0,-3cm) circle (2.5cm) node (current page.center) {\includegraphics[width=5cm]{logo.jpg}};
        \end{tikzpicture}

    }

    \vfill

\end{titlepage}

\restoregeometry






















\newpage

%---------------header------------------%
\pagestyle{fancy}
\fancyhead{} % pulizia degli header
\lhead{%
\begin{tikzpicture}
    \clip (0,0) circle (0.5cm);
    \node at (0,0) {\includegraphics[width=1cm]{logo}};
\end{tikzpicture}%
}
\chead{\vspace{\fill}\Title\vspace{\fill}}
\rhead{\vspace{\fill}\Version\vspace{\fill}}
%quad è una spaziatura
%---------------------------------------%

\begin{table}[!h]
	\caption{Versioni}
	\footnotesize
	\begin{center}
		\begin{tabular}{ l l l l p{6cm} }
			\hline                                                             \\[-2ex]
			Ver. & Data       & Redattore   & Verificatore   & Descrizione     \\
			\\[-2ex] \hline \\[-1.5ex]
			1.0  & 2024-03-16 & Elena Ferro & Leonardo Baldo & Stesura verbale \\
			\\[-1.5ex] \hline
		\end{tabular}
	\end{center}
\end{table}

\newpage

\tableofcontents

\newpage

\section{Dettagli della riunione}


\textbf{Sede della riunione}: Google Meet\\
\textbf{Orario di inizio}: 15:30\\
\textbf{Orario di fine}: 16:30\\

\begin{flushleft}
	\begin{table}[!h]
		\begin{tabular}{ |l|l|l| }
			\hline
			\textbf{Partecipante} & \textbf{Ruolo}  & \textbf{Presenza} \\
			\hline
			Antonio Benetazzo     &                 & Presente          \\
			Davide Malgarise      &                 & Presente          \\
			Elena Ferro           & Redattore       & Presente          \\
			Leonardo Baldo        & Verificatore    & Presente          \\
			Matteo Tiozzo         &                 & Assente           \\
			Raul Seganfreddo      & Amministratore* & Presente          \\
			Valerio Occhinegro    &                 & Presente          \\
			\hline
		\end{tabular}
	\end{table}
	\textbf{Partecipanti esterni}: Andrea Dorigo, Daniele Zorzi, Fabio Pallaro.\\
\end{flushleft}
\* Raul Seganfreddo ha partecipato in qualità di amministratore, in sostituzione di Matteo Tiozzo, che non ha potuto partecipare all'incontro.


\section{Ordine del giorno}
\begin{itemize}
	\item Stato di avanzamento della documentazione
	\item Stato di avanzamento dello sviluppo
	\item Considerazioni su tecnologie alternative a ClickHouse
	\item Obiettivi per il terzo sprint
	\item Decisioni prese
\end{itemize}

\newpage

\section{Verbale}

\subsection{Stato di avanzamento della documentazione}
L'incontro è iniziato con un riassunto da parte del responsabile, Davide Malgarise, di quanto svolto nel terzo sprint.
Per quanto riguarda la documentazione, il gruppo si è occupato di portare a termine l'\textit{Analisi dei Requisiti}
ed effettuare una prima stesura delle metriche che costituiscono il cruscotto di valutazione della qualità, all'interno del \textit{Piano di Qualifica}.
La proponente si ritiene soddisfatta delle metriche proposte e dell'impostazione dei grafici relativi.
I contenuti di tale documento verranno ulteriormente discussi quando il gruppo ne porterà a termine la stesura.

\subsection{Stato di avanzamento dello sviluppo}
Il gruppo ha poi provveduto a mostrare una dashboard realizzata con \textit{Grafana},
la quale visualizza i dati relativi a un sensore salvati su \textit{ClickHouse}.\\
La proponente ha suggerito di raffinare la visualizzazione delle dashboard, in particolare:
\begin{itemize}
	\label{item:consigli}
	\item assicurarsi che nei \textit{marker} rappresentanti i sensori nelle mappe siano presenti più di una cifra decimale per la latitudine e la longitudine;
	\item fornire all'utente la possibilità di filtrare i dati visualizzati in base a determinati parametri (ad esempio il sensore a cui si riferiscono);
	\item rendere più intuitivo il grafico "daily mean temperature", aggiungendo una \textit{label} sull'asse x che rappresenti il giorno (al momento presente solo l'orario);
	\item rimuovere il grafico "average temperature per minute" perché poco significativo. Convertirlo piuttosto in un grafico che mostri la temperatura in tempo reale di un sensore.
\end{itemize}

\newpage
\subsection{Sviluppo di \textit{unit test} e \textit{integration test}}
Il gruppo ha inoltre posto una domanda riguardante la necessità di sviluppare \textit{unit test} e \textit{integration test} per il codice prodotto.
La proponente afferma che per il \textit{Proof of Concept} non è obbligatorio, in quanto l'architettura del sistema è ancora in fase di definizione
e potrebbe subire variazioni; eventuali test potrebbero essere sviluppati su componenti che si ritengono stabili e si prevede di continuare a mantenere.

\subsection{Considerazioni su tecnologie alternative a \textit{ClickHouse}}
Successivamente si è brevemente discusso sulle motivazioni per cui \textit{ClickHouse} è stato proposto da \textit{Sync Lab S.r.L.} per il progetto. La proponente ha spiegato che
rispetto a concorrenti come \textit{InfluxDB} e \textit{TimescaleDB}, \textit{ClickHouse}
consente di effettuare più efficientemente query analitiche complesse.
I punti di forza di \textit{InfluxDB} e \textit{TimescaleDB} sono invece una maggiore efficienza e minore richiesta di risorse, tuttavia
hanno una limitata capacità di effettuare query complesse.

\subsection{Obiettivi per il terzo sprint}
Sono stati definiti gli obiettivi per il terzo sprint:
\begin{itemize}
	\item rifinire le dashboard seguendo i consigli di cui discusso al punto \ref{item:consigli};
	\item aggiungere un'ulteriore tipologia di sensore;
	\item popolare il cruscotto di valutazione della qualità all'interno del \textit{Piano di Qualifica};
	\item completare l'\textit{Analisi dei Requisiti};
	\item redigere la parte di \textit{Piano di Progetto} relativa al terzo sprint.
\end{itemize}

\subsection{Decisioni prese}
Il gruppo si impegnerà a comunicare alla proponente i documenti modificati e i loro contenuti con anticipo di qualche giorno rispetto ai \textit{meeting}
di Stato Avanzamento Lavori, in modo da poterli più efficacemente discutere.\\
Inoltre, da ora in poi gli sprint avranno una durata di una settimana.

\newpage
\begin{table}[b]
	\begin{tabular}{@{}p{.5in}p{4in}@{}}
		Data:  & \hrulefill \\
		       &            \\
		       &            \\
		Firma: & \hrulefill \\
	\end{tabular}
\end{table}

\end{document}